\documentclass[dutch]{../khlslides}
\usepackage{graphicx}
\usepackage{pxfonts}
\usepackage{tikz}
\usepackage{calc}
\usepackage{fourier}

\usetikzlibrary{calc,shadows,decorations.markings}


\title{Propositielogica}
\logo{\includegraphics[height=0.5cm]{../KHL.jpg}}
\institute[KHL]{KHLeuven}

\newcommand{\AND}{\wedge}
\newcommand{\OR}{\vee}
\newcommand{\IMPLIES}{\rightarrow}
\newcommand{\IFF}{\leftrightarrow}
\newcommand{\NOT}{\neg}
\newcommand{\union}{\cup}
\newcommand{\intersect}{\cap}


\pgfkeys{
  /tikz/.cd,
  axis/.style={thin,-latex},
  plot/.style={domain=-5:5,thick}
}

\begin{document}

\maketitle

\begin{frame}
  \frametitle{Algebra\"ische Structuren}
  \structure{Wat is een Algebra\"ische Structuur?}
  \begin{itemize}
    \item Een verzameling $A$
    \item Operaties op elementen van $A$
  \end{itemize}
\end{frame}

\begin{frame}
  \frametitle{Voorbeeld \#1: Re\"ele Getallen}
  \begin{itemize}
    \item Verzameling: $\mathbb{R}$
    \item Elementen: re\"ele getallen
    \item Operaties
          \begin{center}
            \begin{tabular}{lc}
              \bf Naam & \bf Notatie \\
              \toprule
              Optelling & $x + y$ \\
              Aftrekking & $x - y$ \\
              Vermenigvuldiging & $x \cdot y$ \\
              Deling & $x / y$ \\
              Machtsverheffing & $x^y$
            \end{tabular}
          \end{center}
  \end{itemize}
\end{frame}

\begin{frame}
  \frametitle{Voorbeeld \#2: Verzamelingen}
  \begin{itemize}
    \item Verzameling: verzameling der verzamelingen
    \item Elementen: verzamelingen
    \item Operaties
          \begin{center}
            \begin{tabular}{lc}
              \bf Naam & \bf Notatie \\
              \toprule
              Unie & $A \union B$ \\
              Doorsnede & $A \intersect B$ \\
              Verschil & $A - B$ \\
              Cartesisch product & $A \times B$
            \end{tabular}
          \end{center}
  \end{itemize}
\end{frame}

\begin{frame}
  \frametitle{Voorbeeld \#3: Booleanse Algebra}
  \begin{itemize}
    \item Verzameling: $\left\{ \; \mathbf{true}, \mathbf{false} \; \right\}$
    \item Elementen: $\mathbf{true}$ en $\mathbf{false}$
    \item Kortere notatie: $1$ en $0$, respectievelijk
    \item Operaties
          \begin{center}
            \begin{tabular}{lcl}
              \bf Naam & \bf Notatie & \bf Uitspraak \\
              \toprule
              Conjunctie & $p \AND  q$ & $p$ en $q$ \\
              Disjunctie & $p \OR  q$ & $p$ of $q$ \\
              Negatie & $\NOT p$ & niet $p$ \\ 
              Implicatie & $p \IMPLIES  q$ & als $p$ dan $q$ \\
              Equivalentie & $p \IFF  q$ & $p$ als en slechts als $q$
            \end{tabular}
          \end{center}
          \vskip2mm
    \item Booleaanse uitdrukkingen heten proposities
  \end{itemize}
\end{frame}

\begin{frame}
  \frametitle{Waarheidstabellen}
  \begin{center}
    \Large\bf Conjunctie
  \end{center}
  \Large
  \[
    \begin{array}{ccc}
      p & q & p \AND q \\
      \toprule
      0 & 0 & 0 \\
      0 & 1 & 0 \\
      1 & 0 & 0 \\
      1 & 1 & 1 \\
    \end{array}
  \]
\end{frame}

\begin{frame}
  \frametitle{Waarheidstabellen}
  \begin{center}
    \Large\bf Disjunctie
  \end{center}
  \Large
  \[
    \begin{array}{ccc}
      p & q & p \OR q \\
      \toprule
      0 & 0 & 0 \\
      0 & 1 & 1 \\
      1 & 0 & 1 \\
      1 & 1 & 1 \\
    \end{array}
  \]
\end{frame}

\begin{frame}
  \frametitle{Waarheidstabellen}
  \begin{center}
    \Large\bf Negatie
  \end{center}
  \Large
  \[
    \begin{array}{ccc}
      p & \NOT p \\
      \toprule
      0 & 1 \\
      1 & 0 \\
    \end{array}
  \]
\end{frame}

\begin{frame}
  \frametitle{Waarheidstabellen}
  \begin{center}
    \Large\bf Implicatie
  \end{center}
  \Large
  \[
    \begin{array}{ccc}
      p & q & p \IMPLIES q \\
      \toprule
      0 & 0 & 1 \\
      0 & 1 & 1 \\
      1 & 0 & 0 \\
      1 & 1 & 1 \\
    \end{array}
  \]
\end{frame}

\begin{frame}
  \frametitle{Waarheidstabellen}
  \begin{center}
    \Large\bf Equivalentie
  \end{center}
  \Large
  \[
    \begin{array}{ccc}
      p & q & p \IFF q \\
      \toprule
      0 & 0 & 1 \\
      0 & 1 & 0 \\
      1 & 0 & 0 \\
      1 & 1 & 1 \\
    \end{array}
  \]
\end{frame}


\begin{frame}
  \frametitle{Oefening Waarheidstabellen}
  \large
  \[
    p \IMPLIES (\NOT q \IMPLIES p \OR q \OR r)
  \]
  
  \[
    \begin{array}{cccc@{\quad}c@{\quad}c@{\quad}c}
      p & q & r & \NOT q & p \AND q \AND r & \NOT q \IMPLIES p \OR q \OR r & p \IMPLIES (\NOT q \IMPLIES p \OR q \OR r) \\
      \toprule
      0 & 0 & 0 & \only<2->{1} & \only<10->{0} & \only<12->{0} & \only<13->{1} \\
      0 & 0 & 1 & \only<3->{1} & \only<10->{1} & \only<12->{1} & \only<13->{1} \\
      0 & 1 & 0 & \only<4->{0} & \only<10->{1} & \only<11->{1} & \only<13->{1} \\
      0 & 1 & 1 & \only<5->{0} & \only<10->{1} & \only<11->{1} & \only<13->{1} \\
      1 & 0 & 0 & \only<6->{1} & \only<10->{1} & \only<12->{1} & \only<13->{1} \\
      1 & 0 & 1 & \only<7->{1} & \only<10->{1} & \only<12->{1} & \only<13->{1} \\
      1 & 1 & 0 & \only<8->{0} & \only<10->{1} & \only<11->{1} & \only<13->{1} \\
      1 & 1 & 1 & \only<9->{0} & \only<10->{1} & \only<11->{1} & \only<13->{1} \\
    \end{array}
  \]
\end{frame}

\begin{frame}
  \frametitle{Tautologie\"en}
  \structure{Definitie}
  \begin{center}
    Proposities die altijd waar zijn, ongeacht de waarden van de variabelen
  \end{center}
  \vskip4mm
  \structure{Voorbeelden}
  \[
  \begin{array}{r@{\;}c@{\;}l}
      p & \OR & \NOT p \\[2mm]
      p & \IMPLIES & p \\[2mm]
      \NOT\NOT p & \IFF & p \\[2mm]
      p \AND q & \IFF & q \AND p \\[2mm]
    \end{array}
  \]
\end{frame}

\begin{frame}
  \frametitle{Logische Equivalenties}
  \structure{Definitie}
  \begin{center}
    Tautologie van de vorm $P \IFF Q$
  \end{center}
  \vskip4mm
  \structure{Voorbeelden}
  \[
  \begin{array}{r@{\;}c@{\;}l}
      \NOT\NOT p & \IFF & p \\[2mm]
      p \AND q & \IFF & q \AND p \\[2mm]
      p \OR q & \IFF & q \OR p \\[2mm]
      p \OR (q \OR r) & \IFF & (p \OR q) \OR r \\[2mm]
    \end{array}
  \]
\end{frame}

{
  % \newcommand{\highlight}[2][]{
  %   \draw[opacity=.5,#1] ($ (#2.south west) + (-.1,-.1) $) rectangle ($ (#2.north east) + (.1,.1) $);
  % }
  \newcommand<>{\highlight}[2]{
    \alt#3{{\color{#2}#1}}{{#1}}
  }

  \begin{frame}
    \frametitle{Logische Equivalenties}
    \Large
    \begin{overprint}
      \onslide<1-2>
      \[
        \highlight<2>{p}{red} \IMPLIES \highlight<2>{(\NOT q \IMPLIES p \OR q \OR r)}{blue}
      \]

      \onslide<3-4>
      \[
        \NOT \highlight<3>{p}{red} \OR \highlight<3>{( \highlight<4>{\NOT q}{red} \IMPLIES \highlight<4>{p \OR q \OR r}{blue})}{blue}
      \]

      \onslide<5-6>
      \[
        \NOT p \OR (\highlight<5>{q}{red} \OR \highlight<5>{p \OR q \OR r}{blue})
      \]

      \onslide<7-8>
      \[
        (\highlight<8>{p \OR \NOT p}{red}) \OR q \OR q \OR r
      \]

      \onslide<9-10>
      \[
        \highlight<10>{\highlight<9>{1}{blue} \OR q \OR q \OR r}{red}
      \]

      \onslide<11>
      \[
        \highlight<11>{1}{blue}
      \]
    \end{overprint}
    \vskip2cm
    \begin{overprint}
      \onslide<2-3>
      Logische equivalentie:
      \[
        \begin{array}{rcl}
          (\highlight{P}{red} \IMPLIES \highlight{Q}{blue}) & \IFF & (\NOT \highlight{P}{red} \OR \highlight{Q}{blue})
        \end{array}
      \]
      We kiezen
      \[
        \begin{array}{rcl}
          \NODE{$P$}{P4} & = & p \\
          \NODE{$Q$}{Q4} & = & \NOT q \IMPLIES p \OR q \OR r \\
        \end{array}
      \]

      \onslide<4-5>
      Logische equivalentie:
      \[
        \begin{array}{rcl}
          (\highlight{P}{red} \IMPLIES \highlight{Q}{blue}) & \IFF & (\NOT \highlight{P}{red} \OR \highlight{Q}{blue})
        \end{array}
      \]
      We kiezen
      \[
        \begin{array}{rcl}
          P & = & \NOT q \\
          Q & = & p \OR q \OR r \\
        \end{array}
      \]

      \onslide<6-7>
      Logische equivalenties:
      \[
        \begin{array}{rcl}
          (P \OR (Q \OR R) & \IFF & ((P \OR Q) \OR R) \\
          P \OR Q & \IFF & Q \OR P
        \end{array}
      \]

      \onslide<8-9>
      Logische equivalentie:
      \[
        \begin{array}{rcl}
          (\highlight{P \OR \NOT P}{red}) & \IFF & \highlight{1}{blue}
        \end{array}
      \]

      \onslide<10-11>
      Logische equivalentie:
      \[
        \begin{array}{rcl}
          (\highlight{1 \OR P}{red}) & \IFF & \highlight{1}{blue}
        \end{array}
      \]
    \end{overprint}
  \end{frame}
}

\begin{frame}
  \frametitle{Oefening}
  Er werd een moord gepleegd. Sherlock Holmes komt ter plaatse en
  stelt volgende zaken vast:
  \begin{itemize}
    \item Of de kok, of de butler was in de keuken.
    \item De kok was in de keuken of in de eetkamer.
    \item Als de butler een sigaar rookte, dan was hij niet in de keuken.
    \item Als de kok niet in de eetkamer was, dan rookte de butler geen sigaar.
  \end{itemize}
  Welk eenvoudig besluit kan je trekken (gebruik een waarheidstabel)?
\end{frame}

\begin{frame}
  \frametitle{Oefening}
  Er werd een moord gepleegd. Sherlock Holmes komt ter plaatse en
  stelt volgende zaken vast:
  \begin{itemize}
    \item Of de kok, of de butler was in de keuken.
    \item De kok was in de keuken of in de eetkamer.
    \item Als de butler een sigaar rookte, dan was hij niet in de keuken.
    \item Als de kok niet in de eetkamer was, dan rookte de butler geen sigaar.
  \end{itemize}
  Welk eenvoudig besluit kan je trekken (gebruik een waarheidstabel)?
\end{frame}

\begin{frame}
  Bob, Jan en Koen worden verdacht van belastingontduiking. Elk
  van hen getuigt als volgt:
  \begin{itemize}
  \item Bob: ``Jan is schuldig en Koen is onschuldig''.
  \item Jan: ``Als Bob schuldig is, dan is Koen het ook''.
  \item Koen: ``Ik ben onschuldig, maar ten minste \'e\'en van de anderen is schuldig''.
  \end{itemize}
  De elementaire uitspraken die erin voorkomen zijn ``een persoon is schuldig of
  onschuldig''. We spreken af dat de proposities b, j en k respectievelijk betekenen
  ``Bob is schuldig'', ``Jan is schuldig'', en ``Koen is schuldig''.
\end{frame}

\end{document}



%%% Local Variables: 
%%% mode: latex
%%% TeX-master: t
%%% End: 
